\section*{Chapter \ref{ch:g9:periodic-table-first-look} --- The Periodic Table: A First Look}

\begin{solution}{exo:g9:periodic-table-first-look:1}
In order of increasing \omterm{def:g9:inside-the-atom:atomic-number}{atomic number} $Z$.
\end{solution}

\begin{solution}{exo:g9:periodic-table-first-look:2}
Carbon: \omterm{def:g9:periodic-table-first-look:period}{period} 2, \omterm{def:g9:periodic-table-first-look:period}{group} 14. Magnesium: \omterm{def:g9:periodic-table-first-look:period}{period} 3, \omterm{def:g9:periodic-table-first-look:period}{group} 2. Argon:
\omterm{def:g9:periodic-table-first-look:period}{period} 3, \omterm{def:g9:periodic-table-first-look:period}{group} 18.
\end{solution}

\begin{solution}{exo:g9:periodic-table-first-look:3}
\omterm{def:g9:periodic-table-first-look:metal}{Metals}: iron, copper, aluminium. \omterm{def:g9:periodic-table-first-look:metal}{Non-metals}: sulfur, oxygen, chlorine.
\end{solution}

\begin{solution}{exo:g9:periodic-table-first-look:4}
Potassium (\ce{K}) below sodium; bromine (\ce{Br}) below chlorine.
\end{solution}

\begin{solution}{exo:g9:periodic-table-first-look:5}
Calcium (\ce{Ca}, $Z = 20$) and nitrogen (\ce{N}, $Z = 7$).
\end{solution}

\begin{solution}{exo:g9:periodic-table-first-look:6}
It is in the same group as sodium, \omterm{def:g9:periodic-table-first-look:period}{group} 1, and elements of a group
behave alike.
\end{solution}

\begin{solution}{exo:g9:periodic-table-first-look:7}
A \omterm{def:g9:periodic-table-first-look:metal}{metal}: it is shiny, conducts electricity and forms a positive \omterm{def:g9:ions:ion}{ion}.
\end{solution}

\begin{solution}{exo:g9:periodic-table-first-look:8}
To keep elements with similar properties in the same column, even if it
meant leaving a box empty. The gaps were predictions: the missing
elements were later found with the properties he had described.
\end{solution}

\begin{solution}{exo:g9:periodic-table-first-look:9}
Carbon: $\ce{C} = 12$; oxygen: $\ce{O} = 16$. The gaps ``? $= 68$'' (next to
Al $= 27.4$) and ``? $= 70$'' (next to Si $= 28$).
\end{solution}

\begin{solution}{exo:g9:periodic-table-first-look:10}
\omterm{def:g9:periodic-table-first-look:period}{Period} 2: 8 elements (lithium to neon). \omterm{def:g9:periodic-table-first-look:period}{Period} 4: 18 elements (potassium
to krypton).
\end{solution}

\begin{solution}{exo:g9:periodic-table-first-look:11}
\omterm{def:g9:periodic-table-first-look:period}{Group} 18. Where nothing must react: the argon inside light bulbs and
welding gas, helium in balloons (it does not burn).
\end{solution}

\begin{solution}{exo:g9:periodic-table-first-look:12}
$(27.0 + 114.8)/2 = 70.9$, close to 69.7.
\end{solution}

\begin{solution}{pb:g9:periodic-table-first-look:1}
\textbf{1.} $Z = 32$; \omterm{def:g9:periodic-table-first-look:period}{period} 4, \omterm{def:g9:periodic-table-first-look:period}{group} 14.

\textbf{2.} Above: silicon. Below: tin. Left: gallium. Right: arsenic.

\textbf{3.} A metalloid.

\textbf{4.} It is in the same group as silicon, and elements of a group
behave alike.

\textbf{5.} $(28.1 + 118.7)/2 = 73.4$.

\textbf{6.} $(69.7 + 74.9)/2 = 72.3$.

\textbf{7.} The left--right average (72.3): along a \omterm{def:g9:periodic-table-first-look:period}{period} the atomic
weight grows steadily, element by element, while down a group it jumps
by large steps.

\textbf{8.} $72.6 - 70 = 2.6$.

\textbf{9.} $(28.1 + 118.7 + 69.7 + 74.9)/4 = 291.4/4 = 72.85$, that is
72.9.

\textbf{10.} The average, 72.9, is close to both Mendeleev's 72 and
Winkler's 72.3 (and today's 72.6).

\textbf{11.} An element predicted before anyone had seen it was found with
the properties foretold: the table was not just a list, it could make
predictions.

\textbf{12.} About 72.9.
\end{solution}
